\section{权重初始化}

\subsection{为什么初始化很重要}

\begin{figure}[H]
\centering
\includegraphics[width=0.95\textwidth]{slides-images/slide-035.jpg}
\caption{权重初始化太小的后果：激活值逐层衰减到接近零\protect\footnotemark}
\end{figure}
\footnotetext{Slides 第36页。}

权重初始化是训练深度网络的关键步骤。不恰当的初始化会导致训练完全失败。

以一个 6 层全连接网络（每层 4096 维，ReLU 激活）为例，考虑两种极端情况：

\begin{itemize}
    \item \textbf{初始化太小}（$W \sim \mathcal{N}(0, 0.01^2)$）：激活值随层数指数衰减，到最后几层几乎全为零。梯度也几乎为零，网络无法学习。
    \item \textbf{初始化太大}（$W \sim \mathcal{N}(0, 0.05^2)$）：激活值随层数指数增长，最终溢出。梯度也会爆炸。
\end{itemize}

\begin{figure}[H]
\centering
\includegraphics[width=0.95\textwidth]{slides-images/slide-036.jpg}
\caption{权重初始化太大的后果：激活值逐层增长到极大值\protect\footnotemark}
\end{figure}
\footnotetext{Slides 第37页。}

\subsection{Kaiming 初始化}

\begin{figure}[H]
\centering
\includegraphics[width=0.95\textwidth]{slides-images/slide-037.jpg}
\caption{Kaiming 初始化：使用 $\sqrt{2/D_{in}}$ 作为标准差，各层激活值的统计量保持稳定\protect\footnotemark}
\end{figure}
\footnotetext{Slides 第38页。}

Kaiming He 等人提出的初始化方法（He Initialization）：

$$W \sim \mathcal{N}\left(0, \frac{2}{D_{in}}\right)$$

其中 $D_{in}$ 是该层输入的维度。因子 2 是为 ReLU 激活函数量身定制的（ReLU 平均将一半的值置零，所以需要额外的 $\sqrt{2}$ 来补偿）。

\begin{importantbox}{Kaiming 初始化的效果}
使用 Kaiming 初始化后，各层激活值的均值和标准差\textbf{保持大致恒定}——既不会衰减也不会爆炸。这使得：
\begin{itemize}
    \item 前向传播中信号强度稳定
    \item 反向传播中梯度大小稳定
    \item 训练从一开始就在良好的区域进行
\end{itemize}
\end{importantbox}

\begin{knowledgebox}{Xavier 初始化 vs Kaiming 初始化}
\begin{itemize}
    \item \textbf{Xavier 初始化}（Glorot, 2010）：$W \sim \mathcal{N}(0, 1/D_{in})$ 或 $\mathcal{N}(0, 2/(D_{in} + D_{out}))$。适用于 Sigmoid/Tanh 激活。
    \item \textbf{Kaiming 初始化}（He, 2015）：$W \sim \mathcal{N}(0, 2/D_{in})$。适用于 ReLU 激活（因子 2 补偿 ReLU 的零化效应）。
\end{itemize}
在实践中，PyTorch 的大多数层默认使用 Kaiming 初始化。
\end{knowledgebox}

\subsection{本章小结}

权重初始化虽然只影响训练的第一步，但它决定了优化的起点，对训练能否成功有着决定性影响。Kaiming 初始化通过根据层的输入维度和激活函数类型来校准初始权重的标准差，确保信号和梯度在深层网络中稳定传播。

%% ========== Part 8: 训练技巧 ==========

